\chapter{视频感知、跟踪与检索}
\label{chap:vision-video-recognition}

“刚才把钥匙放在哪里”是Ego4D第一视角经历查询所关注的问题类型，需要把问题对应到过去的画面\scite{vision-ego4d2022}

路口的一辆汽车在十帧中出现，究竟应计作一辆还是十辆？人在杯子旁伸手，究竟是拿起还是放下？一张图可以提供外观与空间关系，回答这些问题却还需要把不同时间的观测联系起来。联系的对象也不同：车辆计数需要身份，拿放判断需要顺序，找到某次拿杯的片段则需要起止时间。

\ref{chap:vision-image-recognition}章已经建立单帧识别、定位和匹配的基础。本章把连续观测组织成可用的时间证据，分别讨论对象判断、身份与范围维护、活动发现及查询检索。各项任务可以共享特征与轨迹，但必须保留自己的输出与评价条件。以下汽车和拿杯序列、分数及代价均为教学构造。

\section{视频任务与可用信息}
\label{sec:vision-video-information}

视频用按时间排列的帧$\bm{X}_1,\ldots,\bm{X}_n$表示，帧的拍摄时刻为$\tau_1,\ldots,\tau_n$。等帧率时，相邻时刻差为$1/f$秒，其中$f$是每秒帧数；可变帧率、丢帧或抽样后，应保存实际时间戳。帧编号相差十，只有在时间采样已知时才能换算成持续时间。

\subsection{输出对象与时间范围}
\label{sec:vision-video-outputs}

同一段拿杯视频可以产生多种标注。“有人拿杯子”是范围内的活动标签；每帧的人物框及共同身份组成框轨迹；杯子的逐帧像素区域组成掩码轨迹；从手接触杯子到杯子离开桌面的区间则描述某种约定下的动作范围。训练时只有视频标签，不能把未提供的逐帧边界当作监督；评价时框位置正确，也不能替代身份或活动正确。

\term{目标跟踪}（object tracking）在连续观测中延续选定对象的位置或区域，使后续结果仍指向同一个实例。给定首帧汽车框，即使不知道它的车型，也可以尝试跟踪；自动发现所有车辆则还需要检测及新目标管理。跟踪返回的身份通常是序列内编号，不是现实世界中的车牌或人物姓名。

裁剪片段通常已经把目标活动置于有限范围，未裁剪视频还含等待、背景和多次重复。若训练只见到以动作为中心的短片段，使用时直接对整小时录像分类，输入与输出条件都发生了变化。镜头切换还可能使图像坐标失去连续性；跨镜头重新认出同一辆车需要额外身份依据。视频时间上的位移描述也不能直接当作带米和秒单位的三维运动。

\subsection{观测范围与信息到达条件}
\label{sec:vision-video-availability}

离线与在线描述信息何时可用，局部片段与长时历史描述处理多少信息。这两个维度可以组合。离线算法可以只看五帧，在线算法也可以保存数小时历史；后者在时刻$\tau$只能使用已经到达的帧。若允许等待$\delta$秒后输出，应把新增后帧和等待时间一起记录，具体范围见图\ref{fig:vision-observation-window}。

经历查询要求知道历史中哪些细节值得保存。Ego4D中的相关任务需要返回过去片段，第一视角的遮挡、相机运动和可见范围又会影响证据是否存在。部分样本有其他传感信息，使用时应逐样本确认，不能统一假设视频都有完整音频或视线记录。

保存每帧全部表示会随时间增长。MovieChat用长短期记忆压缩密集视觉信息，区分全局问题与针对某一时刻的问题\scite{vision-moviechat2024} 压缩后若只剩“有人开车”，后来就难以从这句话恢复车牌字符。Flash-VStream进一步面向持续到达的视频维护概括信息与细节记忆，响应成本要连同帧率、历史长度、硬件和提问时刻解释\scite{vision-flash-vstream2025} 这些系统所用的视觉语言连接已在\ref{chap:vision-multimodal-models}章建立；本章在\ref{sec:vision-video-history-retrieval}节展开片段回查，在\ref{sec:vision-video-proactive-response}节讨论何时响应。历史存在过不等于回答时仍能访问它。

\section{对象与属性的时序识别}
\label{sec:vision-video-object-recognition}

一辆车穿过画面时，车身、车牌和车灯的最佳观测可能出现在不同帧。时间识别的收益来自这些互补证据。直接把所有帧平均起来，可能把模糊或错配的观测也加进来；必须先确定取什么、是否属于同一对象，再决定怎样汇总。

\subsection{观测与对象证据的选择}
\label{sec:vision-video-evidence-selection}

若相机以每秒30帧记录，一辆车只在0.3秒内经过，用每秒一帧的等间隔抽样可能完全漏掉它。改变起始相位也会改变是否采到车辆。提高采样率增加覆盖，但同时增加检测和编码开销；可以先用较粗的时间索引找出运动或对象出现的范围，再在这些范围内加密读取。这样的筛选仍要检查漏选，不能只评价已找到的片段。

整帧表示适合利用场景上下文，对象裁剪更集中于车身或服饰。裁剪之前若没有可信的框，识别器可能读到路面或另一辆车；相机运动和出画面还会使固定裁剪失效。保存帧号、框坐标和清晰度，才能在识别出错时查明模型看到了什么。车牌需要的分辨率与车型不同，不能仅以“已采到一帧”判断证据充分。

\subsection{跨帧证据的汇总与表示}
\label{sec:vision-video-aggregation}

同一辆车有多次观测后，可以合并已得到的类别得分，也可以沿轨迹保留随时间变化的属性，还可以保存局部网格特征供更细的状态判断。三种选择保留的信息不同：类别汇总返回范围内的判断，轨迹累积需要可信身份，密集特征还需要任务预测头才能成为深度或活动输出。

\subsubsection{逐帧得分与可靠观测的汇总}
\label{sec:vision-video-weighted-scores}

设已经确认若干裁剪属于同一辆车，第$t$次观测对类别$c$给出得分$s_{t,c}$。这些得分可以是分类器的原始输出，也可以是归一化后的类别概率；具体选择应固定。以下用归一化输出演示算术，但并不假定这些数值已经校准。

平均汇总让每次观测同等参与，最大值保留最强的单次响应，可靠性加权则让清晰、完整的观测有更大影响。对非负可靠性$r_t$，只要$\sum_t r_t>0$，可写为
\begin{equation}
  \bar{s}_c=\frac{\sum_t r_t s_{t,c}}{\sum_t r_t},
  \qquad
  \widehat{c}=\argmax_c\bar{s}_c\text{。}
  \label{eq:vision-video-weighted-scores}
\end{equation}
分母使权重和为1，避免同一组相对权重整体扩大后改变得分尺度。没有可靠观测时应返回未知或重新取帧，不能把分母为零的情形悄然当作一组零分。

考虑清晰、局部遮挡和严重模糊的三帧。它们对“轿车、货车”的输出分别为$(0.80,0.20)$、$(0.60,0.40)$和$(0.05,0.95)$，可靠性设为$1,0.5,0.1$。均值为$(0.4833,0.5167)$，最大值为$(0.80,0.95)$，二者都选择货车。加权轿车分数却为
\[
  \frac{1\times0.80+0.5\times0.60+0.1\times0.05}{1.6}
  =0.690625\text{，}
\]
货车分数为$0.309375$。这解释了为何应限制模糊误判的影响，但并不证明加权一定改善识别：如果可靠性规则把唯一有效帧误判为不可靠，结果也会变差。

最大值还可能由一次异常高分主导；均值则可能稀释一帧中短暂可见的车牌或标志。人工设定清晰度规则与通过标注学习权重是不同方案。学习方案要交代训练目标、标注来源和验证划分，不能在测试结果上反复挑选有利的权重。

\subsubsection{沿对象轨迹积累属性证据}
\label{sec:vision-video-track-evidence}

跟踪器在每帧返回身份及位置后，区域识别器就能为每个身份单独积累证据。只在关联与可见性可信时，把当前得分写入历史。滑动窗口保留最近$k$次观测，窗外观测完全移除；指数衰减则逐步降低旧证据权重，适合无需保存全部分数的情况。

令$\bm{s}_t$为当前类别得分向量，$r_t$同时表达观测与关联可靠性，$0<\rho<1$是衰减因子。维护分子$\bm{a}_t$和总权重$b_t$：
\begin{equation}
  \bm{a}_t=\rho\bm{a}_{t-1}+r_t\bm{s}_t,\qquad
  b_t=\rho b_{t-1}+r_t,\qquad
  \bar{\bm{s}}_t=\bm{a}_t/b_t\text{。}
  \label{eq:vision-video-decay}
\end{equation}
初始二者都为零，仅在$b_t>0$时输出。无新观测可令$r_t=0$；这时分子、分母同比衰减，旧判断不凭空改变，但应单独标出证据已经陈旧。采样间隔不等时，可按实际时间差设置$\rho=\exp(-\Delta\tau/\kappa)$，其中$\kappa$是以秒计的衰减时间尺度。

以车身颜色为例，过去的红色得分为0.8，新一次可靠观测为0.6，令旧总权重为1、$\rho=0.5$、新权重为1，更新后红色分数为$(0.5\times0.8+0.6)/1.5=2/3$。更新的是统计状态，分类器参数并未重新训练。若遮挡后错误接入另一辆红车，历史积累可能反而加强身份错误；\ref{sec:vision-video-tracking}节的关联可信性决定这里能否安全汇总。

\subsubsection{密集时空表示与局部状态预测}
\label{sec:vision-video-dense-state}

一辆车的类别已经确定，并不表示模型还保留了车轮的位置、局部深度或车灯即将发生的变化。类别汇总把许多位置压成少数得分；若希望分别判断这些局部状态，就需要先留下带空间位置和时刻的特征，再让任务预测头读取它们。这里的密集表示是每个时空网格位置都有一个向量，而不是每个位置已经有正确的语义或几何答案。

V-JEPA 2.1把图像或视频切成小块，用可见部分预测完整输入在教师编码器中的特征。与仅监督被遮蔽小块的目标不同，它也约束可见位置的预测，使这些位置不只是替其他位置汇总全局信息\scite{vision-vjepa21-2026} 联合嵌入预测架构（joint-embedding predictive architecture, JEPA）在这里预测的是表示向量：目标是让缺失观测的表示可由其余观测估计，并没有一个把向量直接变成未来像素的解码步骤。

设视频切块后的索引集合为$I$，可见位置集合为$C$、遮蔽位置集合为$M$，满足$I=C\cup M$且两者不相交。每个位置包含空间网格坐标和时间编号。学生编码器$E_\theta$只处理可见小块；预测器$P_\phi$接收可见特征与携带缺失位置的遮蔽词元，输出$\widehat{\bm{h}}_i\in\mathbb{R}^{d}$。教师编码器$E_{\bar{\theta}}$处理同一次增强后的完整输入，给出目标$\bm{h}^{*}_i\in\mathbb{R}^{d}$。这样先固定两端所对应的位置，再比较向量差，避免把某处的车轮特征拿去监督另一处车身。

为了同时约束缺失与可见位置，一条训练样本的基本目标写成
\begin{equation}
 \begin{aligned}
  L_{\mathrm{dense}}
  &=\frac{1}{|M|}\sum_{i\in M}
    \lVert\widehat{\bm{h}}_i-\operatorname{sg}(\bm{h}^{*}_i)\rVert_1
   \\
  &\quad+\frac{1}{|C|}\sum_{i\in C}\lambda_i
    \lVert\widehat{\bm{h}}_i-\operatorname{sg}(\bm{h}^{*}_i)\rVert_1\text{。}
 \end{aligned}
  \label{eq:vision-video-dense-prediction}
\end{equation}
$\operatorname{sg}$表示目标端停止梯度，即本次误差不直接更新教师。抽样要求$M$、$C$均非空；两个分母分别按各组位置数归一化，避免遮蔽比例直接决定哪组占据更多权重。原方法令$\lambda_i=\lambda/\sqrt{d_{\min}(i,M)}$，其中$d_{\min}$以小块网格计量可见位置到最近遮蔽位置的距离，$\lambda$控制总体强度；可见与遮蔽位置不同，使用的距离应为正。靠近缺失区域的位置得到较大权重，是为了加强局部连接，不是根据下游测试答案挑选车轮区域。

原方法还读取三个中间编码层和最终层，把各层归一化特征沿通道拼接，经小网络压回预测器所需维度，再为四个层级分别预测教师目标。在每层应用式\eqref{eq:vision-video-dense-prediction}，得到多层自监督信号。训练从图像和视频抽样；两种输入分别使用二维、三维切块投影，共享主体编码器，来自不同数据工作节点的梯度汇总后更新学生与预测器。教师按$\bar{\theta}\leftarrow\mu\bar{\theta}+(1-\mu)\theta$缓慢跟随学生，$0<\mu<1$。目标生成与防止表示塌缩的条件可参考\BookChapterRef{chap:data-derived-knowledge-acquisition}{卷八的数据知识获取章}；在本节所用形式中，教师并非真实深度或动作标签。

以单层的三个教学位置算一次误差：遮蔽位置的$L_1$残差为2，两个可见位置的残差分别为1和3，$\lambda=0.2$、到最近遮蔽位置的距离分别为1和4。两可见权重为0.2和0.1，故本层目标为$2+(0.2\times1+0.1\times3)/2=2.25$。若教师某个参数旧值为1、更新后的学生值为1.2、$\mu=0.9$，教师新值为$0.9\times1+0.1\times1.2=1.02$。误差通过学生计算图更新预测，教师只按移动平均更新；这个例子没有假设真实深度等于某个特征分量。

每次更新先抽取样本与空间增强、采样遮蔽区域，运行两端编码，累积各层误差，再反传并更新参数。停止采用既定训练日程和留出任务检查，不把误差下降说成每处特征都已正确。论文主训练使用16帧片段，后续较高分辨率阶段使用更长片段；因此使用某个已发布权重时还要核对其采样、分辨率及训练阶段。推理若只需要特征，保留编码器并输入完整观测，不必继续运行遮蔽训练。

局部任务需要自己的监督。论文的深度协议冻结编码器，用最终层网格特征训练线性预测层，再恢复输入空间位置。短时交互预期还同时读取最后一帧的高分辨率图像和截至该帧的短视频，训练模块把视频特征汇总到最后一帧坐标，与图像特征融合，再预测下一交互对象的框、动词及距接触时间。高层动作预期则把截至动作开始前一秒的编码特征与预测器给出的未来表示合并，训练动作、动词和名词分类头。未来帧在训练中可提供目标，不能作为测试输入；预测的未来表示也不是真实未来观测。论文这些迁移协议冻结预训练主体、另训任务模块，不能把自监督预训练统称为无监督完成了下游输出。

我们可以固定一段拿杯视频，同时检查三项教学输出：整段类别是否为拿杯，手附近网格是否给出合理深度，以及只看伸手阶段时是否预期后续接触。若类别正确而手的深度错误，就不能用前者替后者通过；若预期头答对，也没有生成尚未拍到的未来帧。原论文的消融还显示，单独增加可见位置损失可改善密集任务却降低某些分类成绩，多层监督再改变这一取舍。密集监督的收益要按输出和训练配置核验，不能写成所有任务同时获益的保证。

\subsection{属性与状态的连续判定}
\label{sec:vision-video-continuous-attributes}

车身颜色在短录像中相对稳定，车灯是否点亮却会改变。对变化属性，应输出“哪辆车在何时处于什么状态”，不能用整个轨迹的均值替代状态序列。一个可核验的规则是仅在连续若干次可靠观测都超过阈值时改变状态，并在长时间遮挡后回到未知。阈值和持续长度需要依据事件标注与误报、漏报代价选择。

例如灯光分数连续为$0.2,0.8,0.3,0.9,0.9$，设阈值为0.7、需要连续两次支持，则第三次出现的低分打断先前支持，到第五次才确认点亮。规则能过滤短跳变，也会增加确认延迟，并可能漏掉真实的短闪烁。车灯任务若关心的是转向灯频率，就应保留变化而不是强行平滑。时间识别的目标由用途决定，并非输出越稳定越好。

\section{对象状态与空间对应的维护}
\label{sec:vision-video-tracking}

逐帧汽车分类可以判断画面里有什么，却还不知道当前的框是否延续上一帧的同一辆车。跟踪状态因而不仅包含位置，还要保存身份、运动估计、外观或掩码历史，以及最近一次可靠观测的时间。状态维护要能从初始目标出发，解释一次匹配怎样改变后续判断。

\subsection{目标与初始状态的建立}
\label{sec:vision-video-initialization}

单目标任务可以由用户给首帧框、点或掩码，明确接下来延续哪个目标。多目标任务则需由检测器发现对象，给新出现的合格候选建立编号，并维护多个状态。相同的跟踪算法，在人工准确初始化与自动检测初始化下接收到的证据不同，比较时应记录初始化方式。

\subsubsection{SiamRPN++的模板与搜索区域匹配}
\label{sec:vision-video-siamrpn}

给定首帧汽车框，可以保留包含目标的模板$\bm{Z}$，并在下一帧围绕上次位置裁取较大的搜索区域$\bm{S}_t$。SiamRPN++用共享的深层骨干提取两者特征，随后让模板与搜索特征逐通道相关，把“这里像不像目标”变成空间分布的匹配证据\scite{vision-siamrpn2019}

对某一通道$k$，相关操作可概括为
\begin{equation}
  r_{u,v,k}=\sum_{i,j}z_{i,j,k}\,s_{u+i,v+j,k}\text{。}
  \label{eq:vision-video-depthwise-correlation}
\end{equation}
这里$\bm{Z}$与$\bm{S}_t$在公式中指进入相关计算的特征，形状分别为$h_z\times w_z\times d$与$h_s\times w_s\times d$，小写元素$z_{i,j,k}$、$s_{i,j,k}$取出相应位置及通道。$(u,v)$枚举允许模板落入搜索网格的位置，求和沿模板空间范围进行。骨干共享之后，模板与搜索分支还经过各自的卷积投影；每个通道先独立比较，再由后续卷积组合通道信息。分类分支为候选锚框预测目标前景得分，回归分支预测相对锚框的中心及尺度偏移。多层特征各形成预测，再以学习的权重融合，使不同层的信息共同决定候选。

训练样本是同一目标的模板和搜索图像对，以锚框与目标真实框的重叠分配前景、背景和回归监督。深层网络中的填充可能带来位置偏好，原方法通过搜索区域的空间移位采样减轻对目标总在中心的依赖。训练结束后，使用时保留首帧模板，逐帧更新搜索区域和目标框；并不每帧重新训练一个类别分类器。

候选偏移还不是原帧框。设搜索裁剪中的锚框中心和尺寸为$(x_a,y_a,w_a,h_a)$，真实框为$(x^*,y^*,w^*,h^*)$。训练让回归输出接近$\delta_x^*=(x^*-x_a)/w_a$、$\delta_y^*=(y^*-y_a)/h_a$、$\delta_w^*=\log(w^*/w_a)$及$\delta_h^*=\log(h^*/h_a)$。中心差除以锚框尺寸，使不同尺度能使用同一输出规则；尺寸比取对数，使放大和缩小在相反方向上表达。分类交叉熵约束前景判断，正锚框的平滑$L_1$误差约束这四个量：小残差使用二次项，大残差使用线性项，再在抽样图像对上平均并反传到骨干、分支和层间融合权重。原工作从ImageNet骨干参数出发，先训练新增匹配分支，再以较小学习率联合更新骨干；检查点依照既定训练轮次和验证选择，而非跟踪到某帧才停止训练。

推理把预测偏移反解为$x=x_a+w_a\delta_x$、$y=y_a+h_a\delta_y$、$w=w_a\exp\delta_w$、$h=h_a\exp\delta_h$。再把裁剪坐标乘以搜索区域缩放的逆比例并加回裁剪原点，才得到原帧位置。例如教学锚框$(32,32,20,10)$预测$(0.2,0,\log1.2,0)$，得到裁剪框$(36,32,24,10)$；若裁剪原点为$(20,10)$且未缩放，原帧中心就是$(56,42)$。实际选框还要比较窗口与尺度惩罚后的候选，按既定平滑规则更新尺寸，再用更新位置裁取下一帧搜索区。首帧模板仍不变，变化的是位置状态。

假设汽车A在首帧中心$x=40$像素，搜索范围覆盖下一帧的$x=20$至80像素。候选A在$x=48$处得分0.80，另一辆相似汽车B在$x=70$处得分0.85。实际选框还结合尺度变化惩罚和位置窗口，不能直接把最高相关响应叫作正确身份。这些约束表达连续运动的偏好，但当A突然加速、完全遮挡或离开搜索范围时，同样可能把结果拉向错误位置。原方法固定首帧模板，在线改写模板属于额外设计；跨镜头身份恢复也超出了局部搜索的输入条件。

\subsection{当前观测与历史状态的对应}
\label{sec:vision-video-correspondence}

对应可以发生在像素、局部描述或对象框之间。像素对应帮助解释画面怎样移动，对象对应决定哪个观测延续哪个身份。某个像素的位移估计很精确，并不自动把被遮挡前后的两辆相似车区分开。

\subsubsection{像素与局部区域的对应：RAFT}
\label{sec:vision-video-raft}

汽车边缘在相邻帧中移到哪里，可以先按像素位置比较，而不必先决定整辆车的身份。\term{光流}（optical flow）是两幅图像间像素的二维位移场。对第一帧的位置$\bm{p}=(x,y)^\top$，位移$\bm{u}(\bm{p})\in\mathbb{R}^{2}$给出第二帧的取样位置$\bm{p}+\bm{u}(\bm{p})$。横轴向右、纵轴向下，位移以当前网格的像素为单位。这里描述图像中的运动；相机移动、物体运动及透视共同影响它。

RAFT先为两帧提取空间特征，在低分辨率特征网格上计算所有位置对的相关性。设两帧各有$m$个位置、每个位置$d$维，则相关矩阵的元素为
\begin{equation}
  c_{i,j}=\sum_{k=1}^{d}f_{i,k}g_{j,k}\text{。}
  \label{eq:vision-video-all-pairs}
\end{equation}
它为第一帧的每个位置保留第二帧的全局匹配候选。原方法沿第二帧的两个空间维度池化，形成多尺度相关金字塔，而保留第一帧查询位置的分辨率。当前光流估计指定第二帧中的查询中心，再在各尺度的局部邻域读取相关证据\scite{vision-raft2020}

设第$r$轮光流为$\bm{u}^{(r)}$。共享的门控更新模块读取局部相关、第一帧上下文、当前光流及隐状态，预测增量$\Delta\bm{u}^{(r)}$，再更新
\[
  \bm{u}^{(r+1)}=\bm{u}^{(r)}+\Delta\bm{u}^{(r)}\text{。}
\]
相关矩阵是预先构造的，随着估计移动，读取中心改变；隐状态也保留已处理的信息。这与每次重新运行互不联系的匹配器不同。训练对多轮估计施加带衰减权重的光流误差，使用时执行规定的更新轮数，并把粗网格结果上采样到输入分辨率。

更具体地说，两帧特征位于原图约$1/8$分辨率，相关张量可视为$\mathbb{R}^{h\times w\times h\times w}$，查询位置保留前两个维度。第$\ell$层把第二帧位置池化$2^\ell$倍，读取中心也要除以$2^\ell$；中心通常不是整数，因此用双线性插值得到相邻相关值的组合。将各层读取拼接，再与当前光流和上下文组合，门控循环模块决定保留多少旧隐状态、写入多少新状态。门控改变内部特征，不是直接宣告某像素属于车辆A。

设这轮组合输入为$\bm{x}^{(r)}$、旧隐状态为$\bm{h}^{(r)}$，二者是同一粗空间网格上不同通道的特征。更新门$\bm{z}=\sigma(\operatorname{Conv}([\bm{h}^{(r)},\bm{x}^{(r)}]))$决定写入比例，重置门$\bm{g}$由另一卷积及sigmoid计算，决定旧状态参与候选的比例；方括号表示沿通道拼接。候选状态为$\widetilde{\bm{h}}=\tanh(\operatorname{Conv}([\bm{g}\odot\bm{h}^{(r)},\bm{x}^{(r)}]))$，再取$\bm{h}^{(r+1)}=(1-\bm{z})\odot\bm{h}^{(r)}+\bm{z}\odot\widetilde{\bm{h}}$。各门与隐状态形状相同，$\odot$表示逐元素乘法；从新状态预测二维光流增量。门参数在不同更新轮共享，循环状态变化与网络参数训练是不同过程，可参考\BookChapterRef{chap:recurrent-networks}{卷五的循环网络章}。

训练一条样本包含两帧、真实二维位移以及有效监督位置集合$V$。若一共展开$R$轮更新，为使较晚结果占更大权重，可计算
\begin{equation}
  L_{\mathrm{flow}}
  =\sum_{r=1}^{R}\gamma^{R-r}
    \frac{1}{|V|}\sum_{\bm{p}\in V}
    \lVert\widehat{\bm{u}}^{(r)}(\bm{p})-\bm{u}^{*}(\bm{p})\rVert_1\text{，}
  \label{eq:vision-video-raft-loss}
\end{equation}
其中$\widehat{\bm{u}}^{(r)}$已经恢复到原图像素单位，$\bm{u}^{*}$为真实位移，$\gamma$在原实验中取0.8。有效集合来自标注与数据协议，$|V|>0$时才计算该样本监督，不因某处误差大就把它从评价中删除。随机初始化的编码与更新模块先在带光流的合成数据训练，再按目标数据微调；批次中运行完整展开并累积误差，用AdamW更新共享参数。固定训练日程后保留模型，测试的更新轮数另行规定。原方法还对递推中的旧光流分支停止梯度，训练的计算图不能简单理解为无限轮迭代。

上采样既要恢复空间网格，也要恢复位移单位。原方法为每个高分辨率位置学习邻近九个粗网格光流的非负、和为1的权重$\omega_{\bm{p},j}$，取$\widehat{\bm{u}}(\bm{p})=8\sum_{j=1}^{9}\omega_{\bm{p},j}\bm{u}_{j}$。因粗网格一步约对应原图八像素，前面的8不能漏掉；掩码权重经softmax归一化，使输出是局部粗位移的凸组合，而不是重新做一次全局匹配。

例如汽车边缘实际向右移动5个像素，教学更新可以从0开始，依次增加3、1.5和0.5，得到$0,3,4.5,5$。这串数值只演示增量与读取位置的关系，不是RAFT收敛的保证。全相关存储随位置数按$m^2$增长，输入分辨率必须计入资源预算；遮挡位置没有可见的第二帧对应，网络给出的稠密值仍是估计。

要恢复三维速度，还须有深度、相机和时间标定。VGGT-$\Omega$把前馈重建扩展到静态与动态场景，可作为从二维对应转向动态空间恢复的延伸入口\scite{vision-vggt-omega2026} 联合预测仍须说明坐标、尺度和可见性，不能把图像位移直接换成真实世界速度；\ref{sec:vision-video-d4rt}节将具体区分点状态时刻和相机参考系，也不能据单个新模型的局部实验推断所有动态场景都已解决。

\subsubsection{对象与已有轨迹的对应}
\label{sec:vision-video-association}

对象级关联把当前检测框与历史轨迹配对。每条轨迹先预测当前可能位置，再比较候选检测；同一检测不能同时延续两条轨迹，同一轨迹在当前帧也不应接收两个独立目标。因此，逐个选择最相似框并不足够，还要约束配对之间的冲突。

\paragraph{运动预测与一对一分配}
\label{sec:vision-video-motion-assignment}

可用框中心、宽高及它们的速度组成状态$\bm{x}_t$，检测框提供位置与尺寸观测$\bm{z}_t$。线性运动近似下，预测与观测关系为
\[
  \bm{x}_{t|t-1}=\bm{F}\bm{x}_{t-1|t-1},
  \qquad
  \bm{z}_t=\bm{H}\bm{x}_t+\bm{\varepsilon}_t\text{。}
\]
$\bm{F}$按帧间时间推进位置，$\bm{H}$取出框中可观测的量。状态含噪声时，卡尔曼更新利用预测与检测的不确定性调整修正量；此处的作用是把不确定的历史预测和当前测量合并，而非宣布运动一定匀速。

以中心、宽高及相应速度组成八维教学状态时，$\bm{H}\in\mathbb{R}^{4\times8}$只读取前四项，$\bm{F}\in\mathbb{R}^{8\times8}$将速度乘时间差加到位置与尺寸。ByteTrack采用的具体尺寸参数化可以不同，不能把这套教学状态当作它的逐字段实现。设预测协方差为$\bm{P}^{-}$、观测噪声协方差为$\bm{R}$，匹配确定后用创新$\bm{e}=\bm{z}-\bm{H}\bm{x}^{-}$计算
\[
  \bm{K}=\bm{P}^{-}\bm{H}^{\top}
  (\bm{H}\bm{P}^{-}\bm{H}^{\top}+\bm{R})^{-1},
  \qquad
  \bm{x}^{+}=\bm{x}^{-}+\bm{K}\bm{e}\text{。}
\]
逆矩阵要求观测空间的创新协方差可逆，实际可以用线性求解而非显式求逆。过程噪声在预测协方差中累计；观测越不确定，更新通常越依赖预测。没有匹配就只保留预测，并标明最近实测时刻。教学一维中心预测为48、检测为50，预测与检测方差分别为4和1时，增益为$4/(4+1)=0.8$，修正中心为49.6，更新方差为0.8。身份由之前的分配决定，不由这个加权平均重新生成。

上述状态平均同时改变不确定性。若过程噪声协方差为$\bm{Q}$，预测为$\bm{P}^{-}=\bm{F}\bm{P}^{+}_{\mathrm{old}}\bm{F}^{\top}+\bm{Q}$，表示历史不确定性经运动传播后再加入新的过程噪声。观测更新可用$\bm{P}^{+}=(\bm{I}-\bm{K}\bm{H})\bm{P}^{-}(\bm{I}-\bm{K}\bm{H})^{\top}+\bm{K}\bm{R}\bm{K}^{\top}$，其中$\bm{I}$为八维单位矩阵；一维教学数值就是$(1-0.8)^2\times4+0.8^2\times1=0.8$。框尺寸要在有效正值范围内，持续外推若产生不合理尺寸需要按协议处理；噪声参数由设定或验证确定，并非本次匹配自动学得。

设预测框为$\widehat{\bm{b}}_i$、检测框为$\bm{d}_j$，可取代价$c_{i,j}=1-\operatorname{IoU}(\widehat{\bm{b}}_i,\bm{d}_j)$，排除重叠太小或运动偏离太大的组合。匈牙利算法（Hungarian algorithm）在允许配对中求总成本最小的一对一分配，即共同选择哪些轨迹接收哪些检测，避免相同检测被重复占用。实际还要允许未匹配，通过虚拟节点或明确的未匹配代价表达，不能强迫每条轨迹接受一个坏检测。

表\ref{tab:vision-association-cost}给出教学代价。若按行先处理轨迹A，它会占用检测1，B只能接检测2，总成本为$0.10+0.80=0.90$。共同分配则选择A接检测2、B接检测1，总成本为$0.20+0.11=0.31$。检测3在空间门限之外，留下供新目标规则检查。全局分配改善的是这个既定代价目标，代价若不能区分真实身份，最优配对也可能错误。
\begin{table}[htbp]
  \centering
  \begin{tabular}{lccc}
    \toprule
    历史轨迹 & 检测1 & 检测2 & 检测3\\
    \midrule
    A & 0.10 & 0.20 & 不允许\\
    B & 0.11 & 0.80 & 不允许\\
    \bottomrule
  \end{tabular}
  \caption{两条轨迹与三个检测的教学关联代价，越小越合适。“不允许”由配对前的门限排除。例中两个有效匹配均优于约定的未匹配选择，检测3不参加这次分配。}
  \label{tab:vision-association-cost}
\end{table}

外观距离可以加入代价，但需固定表示的训练数据、归一化、权重及门限。强相机运动会使所有框一起移动，急转弯会破坏简单速度外推，两辆相似车交换位置则会同时考验运动和外观。应查看错误来自预测、检测还是关联，而不是把整条轨迹失误都归于检测器。

\paragraph{ByteTrack的高低分检测两次关联}
\label{sec:vision-video-bytetrack}

遮挡后的目标可能仍有一个位置合理但得分较低的检测。若检测后立即把它丢掉，关联阶段就失去这次观测。ByteTrack保留这类候选，将得分高低与是否具有可信历史联系起来\scite{vision-bytetrack2022}

原流程把检测分成高分组与低分组，先将高分组与已有轨迹关联，再用低分组匹配第一次未匹配的可继续跟踪轨迹。第二次主要依据运动预测后的框重叠，避免不可靠外观主导遮挡恢复。未匹配的低分框直接舍弃。作者实现还设置低分下限，只让尚处于跟踪状态的未匹配轨迹参与第二次关联；未匹配的高分框经过未确认轨迹处理及新轨迹阈值，才可能建立身份。已丢失轨迹的保留和重激活另有状态与时间限制，不能把所有历史身份无限制地送入低分关联。

设高分门限为0.6，汽车A、B的当前检测得分分别为0.9和0.3。第一次只匹配A，B暂时未匹配；若低分B框与其预测框重叠充分，第二次就能延续B。同时另一个得分0.2的背景框没有合适历史，会被舍弃。图\ref{fig:vision-bytetrack-association}展示这种条件关系。低分并不等于背景，也不等于一定存在目标，历史用于筛选它的用途。

\begin{figure}[htbp]
  \centering
  % 构造示例：横向位置仅示意关联，不表示实测轨迹。
\begin{tikzpicture}[x=1cm,y=1cm,font=\figurefont\small,
  arr/.style={-{Stealth[length=2mm]},BookInk,line width=0.8pt}]
  \node[anchor=west] at (0,3.3) {上一帧};
  \node[anchor=west] at (3.5,3.3) {当前预测与检测};
  \node[anchor=west] at (7.3,3.3) {关联结果};
  \foreach \y/\name/\col in {2.3/A/BookBlue,0.9/B/BookTeal}{
    \draw[\col,line width=1pt] (0.2,\y-0.35) rectangle (1.6,\y+0.35);
    \node at (0.9,\y) {\name};
    \draw[arr] (1.8,\y) -- (3.4,\y);
    \draw[\col,dashed,line width=0.9pt] (3.6,\y-0.45) rectangle (5.1,\y+0.45);
    \draw[\col,line width=1.1pt] (3.8,\y-0.35) rectangle (5.2,\y+0.35);
    \node at (4.5,\y) {\name};
    \draw[arr] (5.4,\y) -- (7.1,\y);
  }
  \node[above] at (6.2,2.3) {高分0.9};
  \node[above] at (6.2,0.9) {低分0.3};
  \node[anchor=west] at (7.3,2.3) {第一次延续A};
  \node[anchor=west] at (7.3,0.9) {第二次延续B};
  \draw[BookMuted,line width=0.9pt] (3.8,-0.8) rectangle (5.2,-0.1);
  \node at (4.5,-0.45) {背景};
  \draw[arr] (5.4,-0.45) -- (7.1,-0.45);
  \node[above] at (6.2,-0.45) {低分0.2};
  \node[anchor=west] at (7.3,-0.45) {无匹配，舍弃};
\end{tikzpicture}

  \caption{两次关联的教学序列。第一次用高分框延续A，第二次只用位置相符的低分框延续B。虚线框表示预测位置，实线框表示当前检测；右侧孤立低分框没有可信历史而被舍弃，不据此新建身份。}
  \label{fig:vision-bytetrack-association}
\end{figure}

检测器的网络训练与使用时的两次分配是两项过程。关联改变框状态和身份记录，没有在线重新训练检测网络。门限与最大丢失时间需要用验证集确定；如果目标完全没有检测候选，低分阶段也没有新证据可用。持续预测可以等它重新出现，却不能作为已经观测到目标的证明。

在这条教学序列中可把第二次关联也算清楚。设B的预测框与低分B框的IoU为0.7，与背景框为0.05，第二次允许IoU至少0.5，则代价0.3的前者可匹配，后者在求分配前被排除。B用检测更新位置、重置丢失计时并继续保留原编号；A已匹配，不再参加第二次分配。另一个高分新框若没有既有对应，还要经过未确认状态与新建门限，不能把它混入B的恢复。若下一帧连低分B也消失，B转为丢失状态，后续只有合格重激活才能把预测重新变为有观测支持的轨迹。

\subsubsection{动态点查询与相机运动的共同解释}
\label{sec:vision-video-d4rt}

移动相机拍摄行走者时，手臂在图像中向右，可能是手臂移动，也可能是相机转向。RAFT描述两幅图像位置的差，对象编号描述实例连续性，两者都没有单独给出共同三维坐标中的运动。D4RT把完整视频编码一次，再按需要查询某个表面点在不同时刻、不同相机参考系中的三维位置，使深度、动态对应和相机关系可以从同一接口恢复\scite{vision-d4rt-query2025}

设视频编码得到$\bm{F}=E_\theta(\bm{V})$。一条查询为$q=(u,v,t_{\mathrm{src}},t_{\mathrm{tgt}},t_{\mathrm{cam}})$：$(u,v)$是来源帧上归一化像素位置，$t_{\mathrm{src}}$指定在哪帧选点，$t_{\mathrm{tgt}}$指定询问这个点在哪个时刻的位置，$t_{\mathrm{cam}}$指定用哪帧相机坐标表示。解码器返回$\bm{p}(q)=D_\phi(q,\bm{F})\in\mathbb{R}^{3}$以及可见性等辅助输出。选点时刻、点状态时刻和坐标系时刻是三项独立条件，不能只用一个“当前帧”省略它们。

查询词元把二维坐标的傅里叶特征、三个离散时间嵌入及来源点周围$9\times9$的RGB小块编码相加。小块提供所选表面点的局部外观，时间嵌入说明要查询的状态与参考系。每条查询独立交叉读取共享视频特征，没有查询间自注意力；因此增减另一批查询不要求当前点先看到其他查询的输出。共享编码器仍利用输入范围的跨帧信息，这种独立解码不等于严格在线。

固定来源点和$t_{\mathrm{cam}}$，改变$t_{\mathrm{tgt}}$，得到共同坐标系的运动；令$t_{\mathrm{cam}}=t_{\mathrm{tgt}}$，则每个输出采用当帧相机坐标，需要相机关系才能相减。三个时间都等于同一帧时，$\bm{p}=(p_x,p_y,p_z)^\top$的第三项提供该像素深度。在主点为$(0.5,0.5)$的归一化针孔条件下，投影满足$u=0.5+f_xp_x/p_z$、$v=0.5+f_yp_y/p_z$，所以可从多个非退化点算$f_x=p_z(u-0.5)/p_x$和$f_y=p_z(v-0.5)/p_y$，再取中位数。靠近分母为零的位置不能直接参与；畸变相机和不同主点需要改变或细化模型。

相机外参也由同一物理状态在两个参考系的点集恢复。固定来源与目标时刻，只把$t_{\mathrm{cam}}$分别设为$i$和$j$，得到两组三维点，求最贴合的刚性关系$\bm{p}^{(j)}\approx\bm{R}_{j,i}\bm{p}^{(i)}+\bm{t}_{j,i}$。点集去中心后用$3\times3$奇异值分解计算旋转，并用中心位置恢复平移；点几何退化或两组尺度不一致时需额外检查。单目结果的尺度依赖监督与评价标定，不能未经核实就把距离记为米。

把两组$N$个对应点分别记作$\bm{a}_k,\bm{b}_k\in\mathbb{R}^3$，先求均值$\bar{\bm{a}},\bar{\bm{b}}$，再形成$\bm{H}=\sum_k(\bm{a}_k-\bar{\bm{a}})(\bm{b}_k-\bar{\bm{b}})^\top$。若$\bm{H}=\bm{U}\bm{\Sigma}\bm{V}^{\top}$，旋转可取$\bm{R}=\bm{V}\operatorname{diag}(1,1,\det(\bm{V}\bm{U}^{\top}))\bm{U}^{\top}$，平移为$\bm{t}=\bar{\bm{b}}-\bm{R}\bar{\bm{a}}$。去均值先消去平移，奇异值分解在对应方向中寻找总平方误差较小的旋转；行列式修正确保旋转没有变成镜像反射。这里按共同尺度恢复刚性变换；若另估尺度，求解和结果含义也要相应改为相似变换。

训练一条样本由视频和一批可获得几何监督的查询组成，目标包括三维点、二维位置、法向、可见性或点运动；某项标注缺失时不计算该项误差。主要三维误差先让预测点集、真实点集分别除以各自平均深度，再逐坐标施加$g(x)=\operatorname{sign}(x)\log(1+|x|)$，比较$L_1$差。深度归一化减轻尺度差，带符号对数减轻远点的数值支配；二者都改变了误差口径，不能把训练目标称作原始米制距离误差。辅助项分别比较二维坐标、法向余弦、可见性二元交叉熵和运动向量，并对三维误差加入学习置信权重及防止置信全降为零的惩罚。

每批只抽取部分有监督的查询，运行共享编码与独立解码，合并有效误差并端到端更新参数。原工作采用VideoMAE预训练编码器，主训练用48帧、$256\times256$视频，每样本解码2048条随机及区域偏置抽样的查询，用AdamW训练50万步；这份配置说明采样和停止依据，不提供收敛到真实场景的保证。使用时参数固定，同一个编码结果可以回答新的点与时间查询；这不同于RAFT用学习的循环模块多次修正同一光流，也不同于逐场景重新优化一套三维参数。密集查询还可把已被可见轨迹覆盖的位置标为访问过，只从未访问位置发起新轨迹，减少重复解码；遮挡处不能仅因模型输出一个点就标成已看到。

我们用同一个手腕点作教学检查。固定参考系0，时刻0、1、2的点分别设为$(0,0,2)$、$(0.1,0,2)$、$(0.2,0,2)$，则每步横向位移为0.1个共同尺度单位。若时刻1相机向右移0.05而方向不变，其相机坐标横向值为$0.1-0.05=0.05$，对应归一化投影为$0.5+1\times0.05/2=0.525$。只相减各帧相机坐标会漏掉相机运动贡献。相机估计有误时，三维轨迹可以平滑却不准确；手腕遮挡后的点需单独检查可见性，表面点编号也不能替代整个人的永久身份。

\subsection{位置、身份与区域状态的更新}
\label{sec:vision-video-state-update}

框轨迹只给出包围区域，抠取汽车、替换背景或计算精细轮廓还需要逐帧掩码。视频目标分割因此把已知目标的像素信息也保存在历史中，用当前图像查询历史的外观与区域，再恢复当前目标范围。

\subsubsection{XMem的键值读取与分层记忆巩固}
\label{sec:vision-video-xmem}

给定首帧图像和目标掩码，XMem用图像特征构成用于比较的键，用图像与目标掩码共同编码对象值；当前图像的键作为查询。匹配历史位置后，模型把对应的对象值加权读出，结合短时隐状态解码当前掩码\scite{vision-xmem2022}

为看清读取含义，设当前查询位置为$i$、历史位置为$j$，匹配分数为$s_{i,j}$，对象值向量为$\bm{v}_j$。归一化读取写为
\begin{equation}
  a_{i,j}=\frac{\exp s_{i,j}}{\sum_\ell\exp s_{i,\ell}},
  \qquad
  \bm{r}_i=\sum_j a_{i,j}\bm{v}_j\text{。}
  \label{eq:vision-video-memory-read}
\end{equation}
原方法用查询的通道选择权重$e_{i,c}\in[0,1]$和记忆项的缩放系数$\alpha_j\geq1$调节键距离：
\[
  s_{i,j}=-\alpha_j\sum_c e_{i,c}(k_{j,c}-q_{i,c})^2\text{。}
\]
其中$\bm{k}_j,\bm{q}_i\in\mathbb{R}^{d_k}$分别是历史键和当前查询，$\bm{v}_j,\bm{r}_i\in\mathbb{R}^{d_v}$分别是对象值与读出值。通道选择决定这次比较关注哪些外观维度，缩放系数改变该记忆项在距离增大时衰减多快，两者均由学习模块产生。实际读取还可只保留最高的若干匹配项。若两个候选读取权重为0.8和0.2，值中某个分量分别为0.9和0.1，则该分量读出0.74，再由解码器与当前特征共同形成像素判断，不能直接把0.74当作最终掩码概率。

三类记忆承担不同时间尺度的职责。感知记忆是快速更新的隐状态，保存近期变化；工作记忆保存近期较完整的键值信息；长期记忆用较少原型保存反复有用的历史。工作记忆达到容量后，方法根据读取使用情况选择原型键，并把候选值按与原型的联系汇聚成长期值。长期容量也需管理，以使用频率等信息淘汰部分内容。

这次巩固并不是把所有旧帧平均成一幅图。设工作记忆最多保存$T_{\max}$个记忆帧，达到容量时仍保留首帧标注和最近$T_{\min}-1$个帧，把其余帧中的键值位置作为候选。候选$j$的使用量是它历次读取权重的累计和除以实际在库时长，按该量选择$P$个原型键；除以时长使较早入库本身不会成为被选中的唯一原因。每个原型再查询所有候选，取
\begin{equation}
  \bm{v}^{\mathrm{p}}_k
  =\sum_j
   \frac{\exp s(\bm{k}^{\mathrm{c}}_j,\bm{k}^{\mathrm{p}}_k)}
        {\sum_\ell\exp s(\bm{k}^{\mathrm{c}}_\ell,\bm{k}^{\mathrm{p}}_k)}
   \bm{v}^{\mathrm{c}}_j\text{。}
  \label{eq:vision-video-xmem-consolidation}
\end{equation}
其中上标$\mathrm{c}$表示候选，$\mathrm{p}$表示原型；键属于$\mathbb{R}^{d_k}$，对象值属于$\mathbb{R}^{d_v}$，相似度沿用前述学习距离。选择键控制保留多少位置，汇聚值则让未被选为原型的候选仍能贡献信息。原型键值加入长期记忆后才删除相应工作候选；长期达到自己的容量再按使用量淘汰。

教学中三个车身候选的使用量设为0.6、0.3、0.1，只保留第一个原型。若该原型查询候选所得权重为0.6、0.3、0.1，对应某值分量为0.9、0.7、0.1，长期值为$0.54+0.21+0.01=0.76$，并非直接复制首项0.9。后来遮挡解除，当前查询同时读取近期值和长期原型；只有长期键仍与真实车身相符，这个0.76才可能提供有用信息。被压掉的原始车牌像素不能由这个分量复原。

学习这套读取需要逐帧分割目标。一条样本从标注视频抽取若干有时间间隔的帧，以首帧掩码建库，依次读取、解码、把预测作为后续条件，最后比较各目标帧掩码。原工作组合等权的自举交叉熵（bootstrapped cross entropy）与Dice损失（Dice loss）。前者先为像素错误计分，再聚焦其中较难的像素；后者用预测与真实区域的重叠约束整体范围。以一个二元目标的像素概率$p_i$和标签$m_i\in\{0,1\}$为例，Dice项可写为$1-(2\sum_i p_im_i+\epsilon)/(\sum_i p_i+\sum_i m_i+\epsilon)$，$\epsilon>0$处理空区域；实际多目标还需按对象与有效像素协议组织。原工作先从变形静态图像生成三帧合成序列预训练，再从标注视频抽取八帧主训练，逐渐扩大采样间隔；每次训练查询至多随机保留三个历史帧，学习隐状态对较长序列的作用。每批误差反传到编码、选择权重、隐状态与解码模块，用AdamW按固定日程更新并保留权重。长时巩固的容量管理是推理规则，不需要给每个长期原型重新求网络梯度。

汽车被遮挡几秒后重新出现，最近一帧可能只剩遮挡物，长期记忆仍可能保存此前清晰车身。读取结果取决于当前键能否再次指向正确原型，而不是只取决于记忆保存得久。训练通过标注掩码学习编码、读取与解码；使用阶段写入的是特征、预测掩码和隐状态，网络参数保持固定。错误掩码写入会污染后续读取，原型压缩也可能丢掉车牌等细节，有限容量不提供长期身份正确的保证。

\subsubsection{SAM 2的提示分割与流式记忆更新}
\label{sec:vision-video-sam2}

若首帧掩码难以一次标好，可以从点或框提示开始，再针对错误追加提示。SAM 2把这种交互带入视频。本节依据论文2024年10月的修订版，其中更新后的模型对应公开实现中的SAM 2.1。当前图像编码特征通过记忆注意力读取历史空间信息和对象指针；对象指针是从目标解码表示形成的紧凑对象信息，用来参与跨帧条件，而不是现实身份标签\scite{vision-sam2}

提示编码器把点、框或掩码转换为条件，掩码解码器把条件与读过历史的图像特征结合，预测目标区域、质量信息和对象存在得分。不存在或被遮挡时仍要维护状态，但不应据一个看似完整的框认定对象当前可见。预测掩码与图像特征经记忆编码后写入历史。近期非提示帧只保留有限数量，用户提供条件的帧按另外的选择规则参与读取，避免长视频无界存储。

训练使用图像和视频标注，并模拟初始提示及根据预测误差追加的纠正点击，使网络学会利用交互输入。使用时，用户在汽车上点一下，模型传播掩码；若后续车尾漏分，再在漏分区域追加前景点，就会改变该帧的掩码和之后可读取的记忆。纠正影响多远，还取决于实际重新传播了哪些帧，不能只修改当前帧就宣称全部历史已修复。

当前图像特征先作帧内交流，再交叉读取空间记忆与对象指针，随后与提示一起进入掩码解码器。空间记忆保留网格，指针是从掩码输出词元形成的短向量，两者没有相同的空间分辨率。记忆编码器把预测掩码下采样，与未读取历史的图像特征相加并融合，得到下次读取的空间值；因此错误掩码会改变保存内容，而不是仅在显示端留一个错误轮廓。对单点歧义，模型输出多个可能掩码；若没有新提示消除歧义，则按预测IoU选一个传播，不能同时把多个不相容解释称作已确认目标。

训练用8帧序列模拟这条交互链，至多选择两个提示帧；初始条件可来自真实掩码、真实区域内前景点或包围框。当前预测与真实掩码的差决定纠正点击的候选区域，网络必须在收到点击后重新计算，才能学到“这处遗漏如何被纠正”。主干用预训练图像编码器初始化，分割、质量和存在预测通过联合图像视频训练适配。掩码的焦点损失与Dice损失分别约束难分像素和整体重叠，预测IoU用绝对误差约束，存在输出用交叉熵约束；真实目标不在帧中时，只监督存在头，不要求凭空生成一个目标掩码。每批展开交互传播，累积有效监督并用AdamW更新，训练按既定日程结束；线上写库不是重新训练这些参数。

在四帧教学序列中，第1帧前景点选中汽车，第2帧读取首帧与近期记忆，第3帧车尾漏分。第3帧追加车尾前景点后，解码器重新输出掩码，这张新掩码经记忆编码成为后续条件；重跑第4帧时才能利用它。若第3帧同时误包含另一辆车，还需负点或更明确的框排除，追加一个前景点不必解决全部歧义。恢复掩码到原图时还要逆转编码前的缩放与裁剪，并保留提示坐标的同一变换，避免点击和被修区域错位。

在线前向传播只能读已到达帧，离线交互可从不同提示时刻向指定方向重新处理给定视频。二者需要不同的观测范围和等待记录。多目标使用时仍需区分各对象的提示、输出及状态；用户点选一辆车后的持续分割，与自动发现并命名所有车辆是不同任务。

\subsection{遮挡、丢失与重新出现的处理}
\label{sec:vision-video-recovery}

可靠匹配后，状态有当前观测支持；暂时失去匹配后，位置来自外推，身份来自历史。输出应能区分这两种来源。等待时间内若出现候选，重新核验空间范围、外观和历史一致性；超过协议期限再终止轨迹，必要时将再次出现的对象作为新身份，而不强行连接。

固定模板避免把错误目标写成新模板，却难以适应外观大变；更新模板能追随姿态与光照，也可能在遮挡时吸收背景。掩码记忆有同样的两难。实际管理可以对低可信观测停止写入、保留早期可靠模板，或允许用户纠正，但这些措施应按具体实现验证。跨镜头相似外观不足以证明同一实例，历史已经淘汰的细节也不能靠延长推理重新获得。

\section{活动与事件的判断}
\label{sec:vision-video-activities}

身份把对象在时间上连起来，活动判断还要解释对象之间发生了什么。“手、杯子和桌子”在拿起与放下时都可能出现，区别在运动方向、接触关系和状态顺序。本节讨论没有外部查询时，从预设活动类别中判断发生了什么，以及是否还要确定发生范围。

\subsection{活动证据的形成}
\label{sec:vision-video-activity-evidence}

外观可以提示场景和对象，光流可以提供局部运动，手与杯子的相对变化则提示交互。证据怎样分布在时间上，决定采样和表示需要保留什么。长动作需要较广的覆盖，短暂运动需要足够密集的观测，两种要求并不相同。

\subsubsection{TSN的分段稀疏采样与类别共识}
\label{sec:vision-video-tsn}

只取相邻几帧可能看到“伸手”，却没有看到杯子之后是否离开桌面。时间分段网络（temporal segment network, TSN）把给定视频均分成$K$段，每段选一个短片段，用共享网络分别输出类别得分，再通过共识函数形成视频判断\scite{vision-tsn2016}

令第$k$段采样片段为$\bm{V}_k$、网络输出为$\bm{s}_k=f_\theta(\bm{V}_k)$。采用均值共识时，先汇总得分，再做softmax及视频标签损失：
\begin{equation}
  \bm{g}=\frac{1}{K}\sum_{k=1}^{K}\bm{s}_k,\qquad
  p_c=\frac{\exp g_c}{\sum_j\exp g_j},\qquad
  L=-\log p_y\text{。}
  \label{eq:vision-video-tsn}
\end{equation}
这里$y$是给定范围的真实活动类别。损失从视频汇总结果反传到共享网络，而不是把视频标签机械当作每一帧都正在发生的活动。训练时在各段随机取样，使用时固定抽帧数量、位置和空间裁剪协议。

均值为何能把同一个视频标签用于多段学习，可以从导数看出。softmax交叉熵对汇总得分的导数为$p_c-\bm{1}[c=y]$，而$g_c$对每段$s_{k,c}$的导数为$1/K$，所以
\[
  \frac{\partial L}{\partial\theta}
  =\sum_{k=1}^{K}\sum_c
    \frac{p_c-\bm{1}[c=y]}{K}
    \frac{\partial s_{k,c}}{\partial\theta}\text{。}
\]
这里$\bm{1}[c=y]$是类别一致时为1的标量指示量。各片段共用$\theta$，一次更新把各段贡献相加，而不是分别训练$K$个互不相干的分类器。目标是在视频层面降低错误，并没有要求每段独立分类都等于$y$；把视频标签复制为逐帧真实标签会改变监督含义。

训练从图像预训练骨干开始，每批抽取若干带视频标签的样本，在各段随机采片段并作空间增强，运行共识目标后以小批量随机梯度更新。RGB输入提供外观，光流输入需从相邻帧计算二维位移并堆叠成通道，其首层初始化还要适配通道数。原工作采用预先约定的训练日程与归一化配置，停止后保留各分支参数和融合权重。使用时可重复多组规定采样与裁剪再平均预测，但必须把这些读取次数计入开销；固定数量的稀疏取样仍没有覆盖保证。

在九秒教学视频中，伸手、拿杯和收手分别主要发生在三个时间段。连续抽取开头三帧可能都只见伸手，跨三段取样却有机会覆盖完整过程。图\ref{fig:vision-temporal-sampling}对照了覆盖范围。若三段对“拿杯、放杯”的得分依次为$(1,0)$、$(3,0)$、$(2,0)$，均值为$(2,0)$，拿杯概率约为0.881；交换这三个得分的顺序，均值完全不变。因此，共识汇总本身不编码段间次序。

\begin{figure}[htbp]
  \centering
  % 九秒教学序列，样本线不表示实际方法的固定帧数。
\begin{tikzpicture}[x=0.82cm,y=1cm,font=\figurefont\small]
  \foreach \a/\b/\name in {0/3/伸手,3/6/拿杯,6/9/收手}{
    \fill[BookPaper] (\a,2.45) rectangle (\b,3.05);
    \node at ({(\a+\b)/2},2.75) {\name};
    \draw[BookRule,line width=0.5pt] (\a,0.2) -- (\a,3.1);
  }
  \draw[BookRule,line width=0.5pt] (9,0.2) -- (9,3.1);
  \draw[-{Stealth[length=2mm]},BookInk,line width=0.8pt]
    (0,0) -- (9.3,0);
  \node[anchor=east] at (9.3,-0.65) {时间（秒）};
  \foreach \x in {0,3,6,9}{
    \draw[BookInk,line width=0.7pt] (\x,0.05) -- (\x,-0.05);
    \node[below] at (\x,-0.08) {\x};
  }
  \node[anchor=east] at (-0.2,1.95) {跨段取样};
  \draw[BookMuted,line width=0.6pt] (0,1.95) -- (9,1.95);
  \foreach \x in {1.5,4.5,7.5}{
    \draw[BookBlue,line width=1.5pt] (\x,1.7) -- (\x,2.2);
  }
  \node[anchor=east] at (-0.2,0.95) {连续片段};
  \draw[BookMuted,line width=0.6pt] (0,0.95) -- (9,0.95);
  \fill[BookTeal!12] (3.6,0.65) rectangle (5.4,1.25);
  \foreach \x in {3.7,3.95,4.2,4.45,4.7,4.95,5.2}{
    \draw[BookTeal,line width=1pt] (\x,0.7) -- (\x,1.2);
  }
\end{tikzpicture}

  \caption{稀疏覆盖与连续片段提供的时间证据不同。图中三段和动作位置为教学设定，竖线表示采样帧。跨段采样覆盖伸手、拿杯和收手，但稀疏点间仍可能漏掉短动作；连续片段保留局部运动，却未必覆盖全程。}
  \label{fig:vision-temporal-sampling}
\end{figure}

TSN还可分别以RGB图像和堆叠的光流作为两路输入，各自形成预测后按规定权重融合。光流要由帧间估计取得，其计算成本不应从系统延迟中删除。RGB分支可能依靠杯子和桌子的静态线索，光流分支则多了运动信息；即使融合后活动类别正确，也没有自动给出精确起止边界。

\subsubsection{I3D的连续片段与时空特征}
\label{sec:vision-video-i3d}

要区分一只手正向杯子靠近还是远离，可把连续帧一起输入，让局部滤波同时跨越空间和时间。I3D将图像网络的二维滤波器扩展为三维滤波器，在保留已有图像参数的同时开始学习相邻时刻的关系\scite{vision-i3d2017}

若二维卷积核为$\bm{W}$，把它沿时间复制$q$份，并将每份设为$\bm{W}/q$。对没有运动、每帧相同的输入，远离时间填充边界的卷积响应满足
\begin{equation}
  \sum_{r=1}^{q}\frac{\bm{W}}{q}*\bm{X}
  =\bm{W}*\bm{X}\text{。}
  \label{eq:vision-video-inflation}
\end{equation}
这个初始化使静态重复片段的响应尺度与原图像卷积对应。随后视频训练让不同时间切片的权重分别变化，网络才可能学习先后和运动；“复制”本身尚未学会方向。池化的时间核和步长也要单独设置，不能把所有空间下采样原样照搬到时间维。

若二维核形状为$k_h\times k_w\times d_{\mathrm{in}}\times d_{\mathrm{out}}$，扩展后的核多出长度$q$的时间维。对一般连续输入，响应是$\sum_{r=1}^{q}\bm{W}_r*\bm{X}_{t+r}$；静态重复时每项都看同一图，才可把$\bm{W}_r=\bm{W}/q$的$q$项相加得到原二维响应。教学中二维滤波对静态手部给出6，复制三份后每份贡献2，总和仍是6。若视频训练后某个标量时间核变成$(-1,0,1)$，输入局部亮度$(2,3,5)$得3，倒序$(5,3,2)$得$-3$，展示时间权重如何表达方向；这组手设权重不是I3D实际学得的参数。

一条分类样本包含连续RGB或光流片段及所属视频的活动标签，三维卷积产生$\mathbb{R}^{t'\times h'\times w'\times d'}$特征，再按网络范围汇总为类别得分。以单标签为例，误差为真实类别的负对数概率，在抽样片段上求批均值，反传到已扩展时间核和分类头。训练可以从二维权重膨胀后在Kinetics学习，也可以直接从已训练的I3D权重出发适配；应记录哪一路和哪些参数被更新。按既定学习率与训练日程停止后，保留权重、帧间隔与预测汇总协议，推理无须为每段新视频拟合时间核。

同样数量的输入帧仍可能对应不同观察范围。若输入16帧、原帧率30、每隔2帧取一帧，首末观测间隔是$15\times2/30=1$秒；改变原帧率或步长就会看到不同速度和范围。两路输入的来源和计算预算也应分别注明，光流生成不能从整体成本中省略。

连续片段保留局部顺序证据，但时间下采样、汇总和有限感受范围仍可能丢掉关键变化。拿杯与放杯在局部看似相反，实际镜头运动、遮挡和背景偏置可能使分类器依赖其他线索。检验运动判断时，可以比较正放、倒放及保持背景而替换动作的受控片段，结果须按具体模型确认。

\subsection{活动类别与发生范围的确定}
\label{sec:vision-video-activity-output}

已经取得时间特征后，输出还可能是一个范围的类别、某个人的动作标签、若干事件区间，或者每个时刻的步骤。网络都在处理视频，并不意味着这些输出可以用同一组标注训练或用同一个分数评价。

\subsubsection{给定时间范围内的活动判定}
\label{sec:vision-video-clip-label}

单标签分类假设一个片段选一个主要类别，通常用softmax及类别交叉熵。若同一人同时站立、拿杯并说话，就需要多个独立标签，可用各类别的sigmoid输出及二元交叉熵训练。阈值决定最终保留哪些标签，类别不均衡与标注遗漏都影响阈值选择。

整段标签还没有说明是谁在做。AVA对关键时刻的人物框标注多种原子动作。原论文方法从关键帧产生人物候选，将相同空间区域在时空特征上汇集，再预测该人的动作标签和框修正，结合人物框与动作监督学习\scite{vision-ava2018} 这一路线把人物区域引入活动判断，但用同一空间范围跨时刻取特征，并不等于已经恢复了精确人物轨迹。

教学画面中，左边的人拿杯、右边的人看手机，两人都站立。整帧输出“拿杯、看手机、站立”没有完成动作到人物的绑定；每个人的多标签才表达所需关系。反过来，人物框和轨迹完全正确也不能替代手物交互的判断。AVA关键时刻的标注规定了评价时刻，不直接给出整次拿杯动作从开始到结束的连续范围。

AVA原论文方法从关键帧的ResNet-50特征生成候选框，再在RGB和光流I3D的每个时间切片上按同一个空间框作区域池化。这是固定空间的短管状区域，不是把框随人物真实运动逐帧搬动。两路特征在特征图层融合，区域与时间汇总后输出人物框修正和各动作的sigmoid得分；大幅移动的人可能部分离开这个区域，正是需要检查的输入局限。

对一个人物候选，设真实多标签向量为$\bm{y}\in\{0,1\}^{C}$、预测为$\bm{p}\in(0,1)^C$，动作误差为$-\sum_c[y_c\log p_c+(1-y_c)\log(1-p_c)]$。每类独立计分，让站立与拿杯可同时为真；人物候选及框回归仍使用框标注监督，不能由这项动作误差自动取得人物位置。原工作以ImageNet初始化关键帧骨干、Kinetics初始化I3D，两路联合更新，以验证集确定训练步数。推理读取指定关键时刻的上下文、生成候选、预测标签和框，再用非极大值抑制处理重复候选。若还要形成整段人物动作轨迹，需要额外关联，关键时刻多标签成绩不能代替该关联的身份核验。

\subsubsection{事件区间与连续步骤的确定}
\label{sec:vision-video-temporal-boundaries}

\term{时间动作定位}（temporal action localization）在未裁剪视频中输出预设动作类别及发生区间，例如两次拿杯分别位于第9至17秒和第25至30秒。\term{动作分割}（action segmentation）则为每个时间位置分配步骤类别，再把相邻同类位置组成片段，例如拿杯、倒水和放杯。区间检测允许空白、重叠和重复事件，逐帧单标签分段通常要求在规定步骤词表中为每个位置作出选择；并发活动需另设输出。

\paragraph{ActionFormer的多尺度时间点与边界回归}
\label{sec:vision-video-actionformer}

不同事件持续时间差别很大，固定长度窗口难以同时贴合短拿杯和长等待。ActionFormer把按时间排列的片段特征编码为多尺度序列，以局部时间注意力交换邻近信息，并通过下采样扩大每层单个位置对应的时间范围\scite{vision-actionformer2022}

各尺度位置分别预测类别及到左、右边界的非负距离。设某位置换算到物理时间后为$\tau$，该尺度一步对应$\Delta$秒，网络输出距离$a,b$以尺度步为单位，则候选区间为
\begin{equation}
  [\widehat{\tau}_{\mathrm{s}},\widehat{\tau}_{\mathrm{e}}]
  =[\tau-\Delta a,\ \tau+\Delta b]\text{。}
  \label{eq:vision-video-boundaries}
\end{equation}
若输入特征曾按视频时长重新采样，还需先恢复其原始时间映射。距离以秒给出时可令$\Delta=1$；不能把不同尺度的原始输出不加换算地相减。

例如$\tau=12$秒、左右距离已经换算为3秒和5秒，输出为$[9,17]$秒。若模型实际在每步0.5秒的网格上输出$a=6,b=10$，仍得到同一区间。图\ref{fig:vision-temporal-boundaries}用这个例子说明时间点与左右边界的关系。一个位置预测正类不表示只有该瞬间发生动作，它还要通过两侧距离描述持续范围。

\begin{figure}[htbp]
  \centering
  % 上下两部分是独立的教学时间范围；纵向位置不承载数值。
\begin{tikzpicture}[x=1cm,y=1cm,font=\figurefont\small]
  \node[anchor=west] at (0,3.7) {时间点预测左右边界};
  \draw[-{Stealth[length=2mm]},BookInk,line width=0.8pt]
    (1,2.45) -- (10,2.45) node[below left] {秒};
  \fill[BookTeal!14] (2,2.85) rectangle (8.4,3.2);
  \draw[BookTeal,line width=1pt] (2,2.7) -- (2,3.35);
  \draw[BookTeal,line width=1pt] (8.4,2.7) -- (8.4,3.35);
  \fill[BookInk] (4.4,3.02) circle (2pt);
  \draw[{Stealth[length=1.7mm]}-{Stealth[length=1.7mm]},BookInk,line width=0.8pt]
    (2,3.02) -- (4.4,3.02);
  \draw[{Stealth[length=1.7mm]}-{Stealth[length=1.7mm]},BookInk,line width=0.8pt]
    (4.4,3.02) -- (8.4,3.02);
  \node[above] at (3.2,3.1) {3秒};
  \node[above] at (6.4,3.1) {5秒};
  \foreach \x/\lab in {2/9,4.4/12,8.4/17}{
    \draw[BookInk,line width=0.7pt] (\x,2.4) -- (\x,2.5);
    \node[below] at (\x,2.4) {\lab};
  }
  \node[anchor=west] at (0,1.55) {逐帧标签与片段碎裂};
  \node[anchor=east] at (1.2,0.9) {真实};
  \node[anchor=east] at (1.2,0.2) {预测};
  \foreach \x/\gt/\pr in {1.5/拿/拿,2.8/倒/倒,4.1/倒/拿,5.4/倒/倒,6.7/放/放,8/放/放}{
    \draw[BookRule,fill=BookPaper,line width=0.5pt] (\x,0.62) rectangle ++(1.2,0.55);
    \node at (\x+0.6,0.9) {\gt};
    \draw[BookRule,fill=BookPaper,line width=0.5pt] (\x,-0.08) rectangle ++(1.2,0.55);
    \node at (\x+0.6,0.2) {\pr};
  }
  \draw[BookAmber,line width=1.1pt] (4.1,-0.08) rectangle (5.3,0.47);
\end{tikzpicture}

  \caption{时间点回归区间与逐帧步骤分段的教学对照。上半图中12秒位置预测左右距离，形成9至17秒区间；下半图为另一段六个等长时间格的标签，孤立的“拿”把真实连续的“倒”分成两段。两种输出需要分别检查边界与碎片化。}
  \label{fig:vision-temporal-boundaries}
\end{figure}

训练时按照动作区间、各尺度允许的回归长度及中心区域选择正样本，对类别使用焦点损失，对正位置的区间使用距离交并比回归损失。中心采样使正监督集中在动作中心附近，但它不是推理时需要输入的真实边界。使用时解码各尺度候选，筛选分数，并用Soft-NMS降低高度重叠候选的重复影响。

输入本身是一段按时间排列的$\bm{Z}\in\mathbb{R}^{n\times d}$片段特征，来源编码器及时间步长是外部条件。投影后局部时间注意力只在邻近位置交换信息，下采样每层约减半序列长度，因而层$\ell$的网格一步覆盖更长时间。类别头和回归头在各尺度共享参数，分别输出$C$个sigmoid得分与两个非负距离。尺度允许的回归范围先限定哪些长度能分配到这层，再用真实区间的中心邻域决定正位置；一个位置若落入多个真实区间，还需按训练实现的分配规则选择目标。

大量背景点让普通分类误差容易由负类主导。焦点损失（focal loss）在二元正确标签概率$p_{\mathrm{true}}$的负对数误差之外乘$(1-p_{\mathrm{true}})^\gamma$，减轻已经容易判对样本的贡献；$\gamma$及类别权重由配置确定。区间回归不仅希望左右距离接近，还希望解码后的范围重叠。令$\widehat{I}$为预测区间、$I^*$为真实区间，时间IoU是交集时长除以并集时长。距离IoU误差可写为$1-\operatorname{tIoU}(\widehat{I},I^*)+(c(\widehat{I})-c(I^*))^2/\ell_{\mathrm{cover}}^2$，其中$c$是区间中心、$\ell_{\mathrm{cover}}>0$是最小共同覆盖区间长度。最后一项在重叠之外约束中心错位；\ref{sec:vision-video-output-protocol}节还将具体复算区间重叠。

论文按全部尺度位置累积分类误差，只在正位置累积回归误差，再除以正样本数量，形成每视频目标；实现要为没有正样本的批次设置非零归一化保护。每批对序列裁剪或填充，并屏蔽无效位置，使用Adam及预热日程更新时间编码器和两预测头；提供预提取特征不等于同步微调原视频骨干。训练按规定日程和验证检查结束，测试可输入较长完整特征序列。Soft-NMS逐次选择高分候选，并按与已选区间的重叠降低剩余分数，减少同一次事件被重复报告；它仍可能把真实紧邻重复动作误抑制，需按次数与边界分别检查。

输出类别来自训练词表，不能自动响应任意文字查询。原方法利用给定序列的前后上下文；若事件正在发生，右端点尚未观测到，在线系统必须区分预测的结束位置与已经确认的结束位置。把完整录像的边界成绩搬到即时告警，会改变可用信息条件。

\paragraph{MS-TCN的逐帧预测与多阶段分段修正}
\label{sec:vision-video-mstcn}

对“拿杯—倒水—放杯”的完整过程，逐帧分类可能在倒水中间偶然跳回拿杯，造成过度分段。MS-TCN以多阶段时间卷积修正概率序列：第一阶段读取视频特征，后续阶段读取前一阶段的类别概率，用逐步扩大的时间感受范围判断孤立跳变是否符合上下文\scite{vision-mstcn2019}

每阶段的膨胀时间卷积跳过一定间隔取样，使多层计算覆盖较长范围。原方法使用非因果卷积，同一位置可以读取未来特征，因而多阶段结构不能直接等同于在线处理。训练在各阶段施加逐帧分类损失，并加入截断的相邻对数概率差：
\begin{equation}
  L_{\mathrm{smooth}}
  =\frac{1}{nC}\sum_{t=2}^{n}\sum_{c=1}^{C}
  \min\!\left\{
  \bigl(\log p_{t,c}-\log p_{t-1,c}\bigr)^2,\ \eta^2
  \right\}\text{。}
  \label{eq:vision-video-smoothing}
\end{equation}
这里$n$为序列长度，$C$为类别数，$\eta>0$为截断上限。原文采用上述$nC$归一化；实现需明确序列首项及无效填充位置如何处理，并避免对零概率取对数。原方法计算这项损失时把前一时刻概率当作固定目标，只对当前$p_{t,c}$一侧求梯度。各阶段的分类与平滑损失加权后求和训练整个网络。

在一次教学跳变中，“倒水”概率从0.8降到0.1，对数差绝对值为$\log 8\approx2.079$，未截断平方约为4.324。若示例设$\eta=2$，这一对只贡献4；从0.8到0.7的变化则仅约为0.0178。截断限制大跳变的惩罚，但分类损失仍需要求真实边界两侧取不同标签，平滑不能独自决定哪些跳变应消除。

最终按最后一阶段概率选取各时刻类别，再合并连续同类标签。图中真实标签为“拿、倒、倒、倒、放、放”，预测为“拿、倒、拿、倒、放、放”，六格中只错一格，却多出一个虚假拿杯段和一段碎裂倒水。仅报告帧准确率$5/6$会掩盖步骤结构错误。相反，过强平滑可能删掉真实的短拿杯步骤，仍须检查边界和重复次数。

第一阶段输入为$\mathbb{R}^{n\times d}$的时间特征，先投影通道，再叠加核宽3、膨胀率为$1,2,4,\ldots$的残差时间卷积。第$\ell$层膨胀率$d_\ell=2^\ell$，步长为1且不下采样时，每层新增$2d_\ell$个网格的感受范围，所以$L$层覆盖长度为$1+2\sum_{\ell=0}^{L-1}2^\ell=2^{L+1}-1$。这是网格长度，还要乘特征时间步长才成为秒；填充边界的部分位置没有实际观测。后续阶段输入$\mathbb{R}^{n\times C}$概率而不是原始像素，学习修正前阶段已有的判断，因此多阶段不增加新的真实视频证据。

一条训练样本包含特征序列和逐位置步骤标签。每阶段目标为$L_{\mathrm{cls}}+\lambda L_{\mathrm{smooth}}$，其中$L_{\mathrm{cls}}=-n^{-1}\sum_t\log p_{t,y_t}$，全部阶段求和后在视频批次上平均。预提取特征的编码器通常不在这项训练中更新，新时间卷积从初始化参数学习，按训练日程用Adam更新并保留模型。平滑项对相邻对数差截断后，大于$\eta$的差不再从这项惩罚获得进一步增大的梯度；前时刻停止梯度又限定了本轮只修当前一侧。这些计算选择能缓和碎片化，但不是“所有输出必须平滑”的硬约束，真实短步骤仍由逐帧标签支持。

\subsection{未完成事件的及时判断与主动响应}
\label{sec:vision-video-proactive-response}

人正靠近热水壶时，未来拿起壶的画面尚未到达。完整录像分类可在动作结束后回看，在线助手却可能需要在手伸向危险位置时提醒。应先约定触发事件、允许响应区间和误报代价，再区分已发生但证据不足、尚未完成和当前无需回答这几种状态。模型预测之后可能发生什么，与确认已经看见什么，是不同输出。

SPOT-Bench把提问和后续视频保持连续，评价主动查询、无请求干预以及多个并存查询的响应，而不只在人工暂停的几个时刻问答。其Timeliness-F1同时检查语义、时机和整条时间线的输出频率，AsynKV则作为无需重新训练的记忆管理基线展示如何组织持续响应\scite{vision-spot-timeliness2026} 这不是把某个离线动作模型直接改名为在线模型；未来帧可见性、静默规则和查询状态都成为新的条件。

设事件的可接受响应区间为$[t_s,t_e]$，输出发生于$\tau$。为了在这个区间内给足额信用，同时惩罚过早与过晚，论文定义
\begin{equation}
 \mathcal{T}(\tau)=
 \begin{cases}
  \exp\!\bigl(-(\tau-t_s)^2/(2\sigma_{\mathrm{early}}^2)\bigr),
    &\tau<t_s,\\
  1,&t_s\leq\tau\leq t_e,\\
  \exp\!\bigl(-(\tau-t_e)^2/(2\sigma_{\mathrm{late}}^2)\bigr),
    &\tau>t_e\text{，}
 \end{cases}
 \label{eq:vision-video-timeliness}
\end{equation}
其中两个$\sigma$以秒为单位，原协议令$\sigma_{\mathrm{early}}>\sigma_{\mathrm{late}}$，即同样偏离长度时更严惩迟到。区间由任务和人类时机标注确定，不能看到模型何时回答再把区间移到那里。时间分数高还需语义正确，内容错误的及时提醒不成为正确匹配。

先考虑每个事件只允许一个有效响应的$K=1$情形。按时间处理所有输出，与尚未占用且语义相符的事件槽匹配，重复或无效输出计误报，未填槽计漏报。设匹配次数为$N_{\mathrm{TP}}$，时间加权信用总和为$x_{\mathrm{TP}}=\sum_{\mathrm{matched}}\mathcal{T}(\tau)$，则精确率与召回率分别为$x_{\mathrm{TP}}/(N_{\mathrm{TP}}+N_{\mathrm{FP}})$和$x_{\mathrm{TP}}/(N_{\mathrm{TP}}+N_{\mathrm{FN}})$，再取调和均值得到Timeliness-F1。分母保留整数响应与事件数，分子才是时间加权信用；不能把它当作只在成功输出中平均的延迟。$K>1$允许每槽规定数量的候选，需遵守论文附录的占用和去重规则，分数比较时固定$K$。

教学事件区间为$[5,6]$秒，设提前尺度1秒、迟到尺度0.5秒。6.5秒的正确输出得到$\exp(-0.5)\approx0.607$；若之后7秒重复提醒，已经填满的槽不能再得信用，计一次误报。此时精确率为$0.607/2$，召回率为0.607，Timeliness-F1约为0.404。完全静默会漏掉该槽，一直重复则降低精确率。对“靠近—伸手—拿起”各时刻还应保存当时可见帧，6.5秒的修订不能回填成5秒已经正确预警。

AsynKV用已有Qwen2.5-VL权重，不添加新的训练目标。每个查询保存编号、提出时刻、是否已回应及最近带时刻的答案；后续上下文含这份答案记忆，使模型知道刚才已经响应。近期帧的键值缓存留在GPU，较早帧移到CPU形成有界历史；收到查询后按相似度取回相关帧缓存，放到近期帧之前并重新安排时间位置。这里缓存是语言模型已经算过的内部键值，不是原始视频，也不能从一个缓存项任意恢复未保留的像素。

没有待回应查询的时段可用于非重叠窗口摘要、缓存搬运和压缩。原方法寻找相邻最相似的两帧，丢弃较旧项来控制长期容量，后台CPU维护检索分数；这些操作减少在线关键路径上的工作，也可能删掉之后需要的细节。原配置每秒读一帧、GPU近期16帧、CPU长期至多256帧，每查询取回4帧，并将查询有效期设为120秒。数字是该实验协议，不是视频助手的一般最优设置。更高帧率或繁忙查询下后台任务仍可能竞争资源，要实际记录硬件、排队、漏帧和响应耗时，不能把“在空闲时执行”写成任何负载都没有延迟代价。

\subsection{顺序、重复与状态变化的解释}
\label{sec:vision-video-event-relations}

得到事件区间后，可以把每个结果保存为对象身份、类别、开始时间、结束时间及来源位置。两个事件是否相邻、重叠或重复，应由这些记录计算，而不是只从视频总标签推测。“拿起”与“放下”还要结合杯子相对桌面的状态变化；车辆转弯则要结合轨迹方向及道路参照。

两次拿杯若因区间过宽而合并，次数会少算；一次倒水若被分成三个小段，次数会多算。前一段结束恰好早于后一段开始，只提供时间先后，并不证明前者导致后者。更复杂的跨事件语言解释将在\ref{chap:vision-multimodal-understanding}章中连接到答案证据核验。

\section{视频匹配与片段检索}
\label{sec:vision-video-retrieval}

前面从视频出发发现对象与事件，检索则从查询出发决定什么内容值得返回。同一组帧、轨迹和片段表示都可以进入索引，区别在查询要求及相关性的定义。返回一个相似视频只完成了候选筛选，是否包含同一辆车、所需动作及准确时间范围还要分别核验。

\subsection{查询与匹配要求的确定}
\label{sec:vision-video-query}

视频查视频可能寻找相同事件或相似运动，图像查视频可能寻找相似外观，给定车辆照片则可能要求同一实例。即使三种查询都用向量相似度排序，正例定义也不同。任务协议还要规定返回单位是整段视频、身份轨迹还是时间区间，以及边界可以偏差多少。

文字查询“有人把杯子放到桌上”增加了语言与视频的对应。图文表示对齐见\ref{sec:vision-models-contrastive}节；这里先把文字编码器视为能够输出可比较向量的接口，重点追踪它与哪些帧、哪些范围比较。精确语言指代及查询区间定位将在\ref{chap:vision-multimodal-understanding}章展开。

\subsection{候选视频与片段的筛选}
\label{sec:vision-video-candidates}

每段视频一个向量便于快速搜索，但把位置压缩掉后很难回答某个局部证据来自哪里。局部索引保存更多表示及时间信息，成本较大，却能把命中回溯到实际观测。选哪种粒度，取决于查询是否需要身份、运动或精确时间。

\subsubsection{帧、轨迹与局部片段的视觉匹配}
\label{sec:vision-video-visual-matching}

图像查询可以使用\ref{sec:vision-recognition-retrieval}节的描述向量，与逐帧全图或对象裁剪比较。活动查询可与连续片段的表示比较。索引条目至少保留视频编号、时间范围及必要的区域坐标，表示来源还应注明训练与归一化方式；建立索引并没有重新学习表示。

给定查询向量$\bm{q}$和候选帧向量$\bm{v}_t$，帧级最大相似度$\max_t\langle\bm{q},\bm{v}_t\rangle$强调曾经出现过的强匹配；轨迹均值则积累同一对象的外观证据。前者可能被一次错误匹配抬高，后者依赖身份关联。对运动事件，独立帧得分无法区分同样画面元素的不同顺序，需要使用保存局部时间关系的片段表示。

相似度计算的候选集合来自实际建立的索引，不能搜索尚未编码或已经删除的位置。若采用余弦比较，先将非零向量除以各自$L_2$范数；零向量应排除或按协议标为未知。教学车辆查询$\bm{q}=(1,0)$，两条轨迹区域表示分别为$(0.8,0.6)$与$(0.6,0.8)$，二者已是单位向量，得分为0.8与0.6，前者先返回。若整帧背景表示反而让第二条得分更高，区域回查就能查出排序依据不同。返回时保留命中帧、对象框和窗口原始时间戳，才可继续判断是哪辆车、何时出现。

以汽车图片查录像，整帧可能因为同一片道路和天空而命中，裁剪区域却显示车型不同。沿命中位置回查框与相邻帧，才能检验对象对应。若索引窗口为10秒，排序只说明这个10秒窗口值得检查，不能认定动作恰好持续10秒。长历史可先检索粗摘要，再读取细粒度条目，但在线只能回查实际保存且已经到达的信息。

\subsubsection{CLIP4Clip的帧聚合与视频排序}
\label{sec:vision-video-clip4clip}

车辆照片可以直接匹配车辆区域，文字查询却还需要把“放杯”这样的描述与视频画面连接。CLIP4Clip从已有图文编码器出发，把采样视频帧编码成图像向量，文字查询编码成文本向量，再学习适合视频检索的比较方式\scite{vision-clip4clip2021} 最直接的配置对帧向量求均值，再与文本计算余弦相似度；序列配置先用长短期记忆网络（long short-term memory, LSTM）或带时间位置的Transformer处理帧序列，再汇总。LSTM按顺序维护特征状态，Transformer用时间位置区分各帧；均值方案本身不保留顺序，两种序列模块则可以把时间关系写入表示。

设一个批次含$B$对匹配视频和文本，$\bm{v}_i,\bm{q}_j$为归一化后的表示，$z_{i,j}$为用于训练的相似度分数。若使用温度$\tau>0$，可令$z_{i,j}=\langle\bm{v}_i,\bm{q}_j\rangle/\tau$。两个方向分别在同一矩阵的行、列中识别正确配对：
\begin{equation}
  \begin{aligned}
    L_{\mathrm{v2t}}&=-\frac{1}{B}\sum_i
      \log\frac{\exp z_{i,i}}{\sum_j\exp z_{i,j}},\\
    L_{\mathrm{t2v}}&=-\frac{1}{B}\sum_i
      \log\frac{\exp z_{i,i}}{\sum_j\exp z_{j,i}}\text{。}
  \end{aligned}
  \label{eq:vision-video-retrieval-loss}
\end{equation}
训练组合两个方向的损失。批内其他配对充当负候选，但如果两段视频实际都符合描述，就可能产生假负例；标注的一一对应并非语义上的唯一匹配。原工作比较冻结范围、学习率和不同聚合器，使用某种配置时应明确哪些预训练参数被微调、哪些新参数从视频文本配对中学习。

这个目标将排序学习限制在当前批的候选集合中。以$B=2$的教学分数矩阵$\bm{Z}=\bigl(\begin{smallmatrix}2&0\\1&3\end{smallmatrix}\bigr)$为例，两行正确配对概率都为$\exp2/(\exp2+1)\approx0.881$，所以$L_{\mathrm{v2t}}\approx0.127$。第一列正确概率为$\exp2/(\exp2+\exp1)\approx0.731$，第二列为$\exp3/(1+\exp3)\approx0.953$，故$L_{\mathrm{t2v}}\approx(0.313+0.049)/2=0.181$。两方向之和约0.308，区别来自各自要区分的候选，不是把一次视频检索误差重复写两遍。单行负对数概率对分数的梯度是预测配对概率减去正确配对指示量；按批次平均再除以$B$，随后传入视频聚合器和允许微调的编码层。

训练样本是视频及其文本描述，不必有动作边界。每批按规定时间位置抽帧，完成图像缩放、裁剪和文本分词，以CLIP权重初始化两编码器；新序列聚合器则需要学习自己的参数。均值配置没有新增聚合参数，仍可能微调编码器，不能把“无参数聚合”理解为完全没有训练。原研究固定图文嵌入底层的部分参数、对不同层使用不同学习率，并比较多种微调范围。按配置以Adam更新双向目标，依训练日程与留出检索成绩选择权重，最终同时保留编码器、聚合器、采样规则和得分尺度。

均值配置使用$\bm{v}=\operatorname{norm}(n^{-1}\sum_t\bm{f}_t)$，其中$n$为采样帧数、$\bm{f}_t\in\mathbb{R}^{d}$为该帧表示，$\operatorname{norm}$将非零均值除以其$L_2$范数。时间索引只用于取样，不参与均值运算。序列配置先将帧表示和时间位置输入循环或Transformer模块，再按配置汇总；其时间交换能力可参考\BookChapterRef{chap:transformer}{卷五的Transformer章}，当前任务仍须用视频文本配对训练聚合结果。二者都必须在图库编码和查询使用时保持相同预处理。若图库变更编码器权重，旧向量不能不加核验地继续与新查询向量比较。

使用时可预先编码候选视频，收到查询后计算相似度并排序。教学查询“有人把杯子放到桌上”对三个视频的分数设为0.82、0.70和0.21，依次返回第一、第二、第三个候选。第二个视频若是拿起杯子，较高相似度可能来自相同对象；若第一段有一分钟而放杯只占两秒，全视频平均又可能稀释这一事件。二者分别指向顺序表达与时间粒度问题。

这一流程排序的是事先给定的视频或片段，没有同时回归动作边界。应用若先切成窗口再检索，窗口长度与重叠就是外部条件，不能把窗口端点解释成模型推断的动作起止。候选筛选的表示与精确定位可以组合，评价仍须分别检查。

\subsubsection{长历史的分层检索与原始片段回查}
\label{sec:vision-video-history-retrieval}

查询“两次敲门都发生在开门之前吗”需要三段事件证据，找到第一次敲门还不能决定答案。长视频均匀抽帧可能漏掉短敲门，全视频均值也会稀释它；如果压缩后只留下“有人在门口”，后来便不能判断敲了几次。可用粗概览确定候选范围，再在局部增加帧密度和分辨率，但每个索引都要指回原始时间戳，且原视频必须仍可读取。

TRACE以原始短片段作视觉锚点，逐轮积累证据并检查答案随证据变化是否稳定\scite{vision-trace-anchors2026} 这里的锚点包含时间范围、抽帧方式和空间分辨率，其值是能直接查看的原始帧，不是上一轮文字摘要。方法不重新训练模型，而是用已有视觉语言模型执行候选提议、证据判断与回查控制。语言交互的通用接口沿用\ref{chap:vision-multimodal-models}章；本节解释它实际取得了哪些视频观测。

初始从完整视频均匀取32帧形成概览，记录其原始时刻。时间提议模块结合概览、问题、选项、已经查看的锚点、未探索区间和未满足证据需要，提出候选范围及观察方式。局部提取按该方式读取原视频，保留与问题有关的短锚点；优先级模块再按区分答案或填补缺口的作用安排锚点。这个顺序是证据加入顺序，不应误记成事件发生顺序，事件本身仍按原始时间戳排列。

设排好序的锚点为$a_{(1)},\ldots,a_{(n)}$。为了查明新增画面是否改变回答，形成前缀证据集合
\begin{equation}
  \mathcal{E}_j=\{a_{(1)},\ldots,a_{(j)}\},
  \qquad
  r_j=\operatorname{EA}(\mathcal{E}_j,q,O)\text{，}
  \label{eq:vision-video-trace-prefix}
\end{equation}
其中$q$是问题、$O$是答案选项，$\operatorname{EA}$是证据判断调用，输出答案候选、可回答或冲突状态、已观察的区分事实及剩余缺口。每次调用只接收当前原始证据、同一问题和选项，不继承上一前缀的答案或文字推理；如此比较$r_j$和$r_{j+1}$时，输入差主要是新增锚点，而非让模型顺着自己旧答案继续解释。

独立前缀的代价是反复输入相同画面。若三个锚点分别含4、3、3帧，三次判断依次输入4、7、10帧，共21帧，而不是只支付10帧。初始32帧概览、局部筛选中给模型看的画面以及最终复查也都计入实际视觉输入成本；用于覆盖检查的不同时间戳数与用于成本检查的重复输入帧数是两种量。缓存解码结果可以减少重复视频读取，却不能把重新送进模型的视觉输入计算当作免费。

轨迹协调模块读取这串判断，查找最早的稳定前缀：后来锚点不改变答案、不降低可回答性、不与已记录事实冲突，且没有未满足证据需要。若画面冗余或错误，就删除锚点并重建判断；若候选里看不清动作，就以更高密度或分辨率重读；若仍缺事件，则扩展未探索范围。即使答案编号没有变化，新锚点补上第二次敲门也说明早先证据并不充分，不能仅凭连续两次回答一样就停止。

在教学视频中，第一个锚点显示10秒敲门，第二个显示40秒开门，当前只能确认一次敲门在开门前，第二次仍是缺口。扩大搜索后，第三个显示25秒另一次敲门，三者时间顺序才支持“两次都在开门前”。若动作模糊，把该范围重读后仍得到同一顺序，协调模块可提出停止。停止前再独立调用一次，只给最早稳定前缀的原始片段、问题及选项；答案一致才报告找到稳定前缀，答案不一致就继续处理冲突。若预算先耗尽，只能回放冲突较少的前缀给出带未稳定状态的强制答案，不能报告充分覆盖。

稳定是模型在已经取得证据上的停止判断，不是录像里绝无第三次敲门的证明。评价还需把真正送给模型的源帧时间戳与每个必要事件区间对照，检查每段是否达到规定观测密度，并单独检查答案。此流程假设离线完整原视频可回查；在线只能在已到达且实际保留的历史范围内执行，摘要里不存在的手势或字符也不能凭更多语言推理补回。

\subsection{检索结果的核验}
\label{sec:vision-video-retrieval-check}

视频级检索检查相关性和排序，实例查询还要核验对象是否同一，区间查询要核验起止与重复次数。比如正确视频中有两次拿杯，只返回覆盖两次的长片段，可能满足“找一段含拿杯的视频”，却不满足“找到第二次拿杯的时间”。查询语义改变了成功条件。

精细核验可以沿候选范围重新检测、关联和判断状态，再结合查询要求选择结果。预设活动检测依赖固定类别标注，任意文字定位则还依赖语言对齐与指代监督。前者不能仅因输出了区间，就被称作已解决后者。

\section{时序判断与历史依据的检验}
\label{sec:vision-video-evaluation}

一次交通分析可以给对车型却交换身份，检索命中正确视频却圈错转弯区间。应先固定共同观察协议，再按输出职责检查错误；使用记忆的系统还要进一步区分“有记忆时更好”与“必须取回正确历史才更好”。

\subsection{输出职责与观察协议}
\label{sec:vision-video-output-protocol}

视频划分应尽量以源视频、对象或场景为单位，避免把相邻帧分到训练和测试后形成近重复泄漏。比较还要固定抽帧、分辨率、初始化、未来帧可见范围和人工纠正预算。离线算法能读完整事件，在线算法可能尚未看到结尾，两者的差异应先解释输入，再比较输出。

框与掩码需要检查空间重叠，跟踪还检查身份连续性。身份F1（IDF1）在预测身份和真实身份作全局一对一对应后，按满足身份及位置条件的观测数计算
\begin{equation}
  \operatorname{IDF1}
  =\frac{2\,\operatorname{IDTP}}
  {2\,\operatorname{IDTP}+\operatorname{IDFP}+\operatorname{IDFN}}\text{，}
  \label{eq:vision-video-idf1}
\end{equation}
其中$\operatorname{IDTP}$为身份正确的匹配观测数，其余两项分别是身份误报和漏报。若同一真实目标的十次观测被错误拆成两个不同身份、长度为6和4，即使框都准确，全局一对一对应也只能选择一条预测身份，得到$\operatorname{IDTP}=6$、$\operatorname{IDFP}=4$、$\operatorname{IDFN}=4$，IDF1为0.6。这个简化例子说明逐帧检测正确不代表身份维护成功。实际还应报告身份切换、轨迹中断及不同遮挡长度的表现。

事件区间用\term{时间交并比}（temporal intersection over union, tIoU）衡量范围吻合。对两个正长度区间$[a,b]$和$[c,d]$，其中$a<b$、$c<d$，交集长度不能小于零，并集长度等于两段长度之和减去交集，故
\begin{equation}
  \operatorname{tIoU}
  =\frac{\max\{0,\min(b,d)-\max(a,c)\}}
  {(b-a)+(d-c)-\max\{0,\min(b,d)-\max(a,c)\}}\text{。}
  \label{eq:vision-video-tiou}
\end{equation}
$[9,17]$与$[10,18]$相交7秒、并集9秒，得分为$7/9$。区间检测的平均精度还要求类别一致、按得分排序并处理同一真实事件的重复命中，宜报告多个tIoU门限。逐帧步骤则同时报告帧准确率与片段级匹配F1或编辑距离，避免长步骤的正确帧掩盖短步骤被删或反复切碎。检索的Recall@$k$和平均精度仍以约定的相关候选为依据。

\subsection{对象、状态与事件错误的共同追踪}
\label{sec:vision-video-error-trace}

交通应用可以把这些结果串起来：检测提供汽车范围，关联维持车辆编号，轨迹与道路方向共同支持转弯判断，检索返回需要复核的片段。错误追踪也应沿同一路径进行。天气导致检测置信降低，可能造成漏配；低帧率使运动预测失准，可能引发身份交换；过宽事件区间又可能把两次转弯算成一次。用一个总分类分数难以识别该修哪一环。

对一条错误记录，可以先回看目标是否确实被采到，再核验候选框、关联代价、掩码写库内容和事件左右边界。若车辆框始终正确而编号中途交换，增加分类器训练并不直接修复关联；若身份不变但转弯区间延长，应检查时间特征和端点恢复。在线记录还要保留首次输出、每次修订及其当时可见的帧，避免把后来借助结尾画面修正的答案当作早期成绩。日志应区分当前观测、历史读出和外推结果，以便确定错误信息从哪里进入。

\subsection{记忆内容替换与历史依赖的诊断}
\label{sec:vision-video-memory-specificity}

汽车遮挡后，使用记忆的分割得分高于无记忆，仍不能证明收益来自这辆车的正确历史。改变记忆还可能改变特征尺度、缓存长度或位置编码；即使放入不相关但形状合适的表示，网络也可能获得类似改善。因此需要固定读取时的其余计算，只替换实际消费的值，再看哪些收益随内容丢失。

记忆特异性（memory specificity）描述记忆收益对所取内容的依赖，而不仅是记忆通路存在与否。相关读取替换审计在已有冻结视频模型、WorldMem和SAM 2上实施了不同对照\scite{vision-memory-specificity2026} 它检查特定读取位置的直接影响，无须重新训练，但要求能访问该模型实际消费的记忆接口；移除整个缓存不是同一项干预。

设$S(\bm{v})$为消费记忆值$\bm{v}$时的任务分数，统一令分数越大越好。正确值为$\bm{v}_{\mathrm{corr}}$，去除内容的参考分数为$S_0$；支持无记忆模式时用该模式，否则保留原接口但将值置零。为了把不同基线上的变化分开，定义
\begin{equation}
  B=S(\bm{v}_{\mathrm{corr}})-S_0,\qquad
  R(\bm{v})=\frac{S(\bm{v})-S_0}{B}\quad(B>0)\text{。}
  \label{eq:vision-video-memory-recovery}
\end{equation}
$B$是正确记忆带来的收益，$R$是替换后还保留的收益比例。只在$B>0$时解释这个比例，不能把没有收益的系统硬写成“恢复多少”。$R$不裁剪到$[0,1]$，因为替换可以低于参考或略优于正确记忆；误差型指标若越小越好，应先规定方向一致的分数，不能直接混用。

在目标读取步骤，查询、已选择位置、条目数量、时间及位置编码、模型参数和此前运行历史全部固定，只将消费值换成有效但错误的记忆。同一轨迹中的错误时刻还保留较多场景上下文，其他轨迹或其他对象的记忆同时改变更广的内容；无特定身份的原型进一步去除实例信息但保持部分表示统计，强扰动则检查偏离训练表示的后果。这些对照移除了不同信息，结果要逐项报告，不能由一种随机噪声就断言系统完全不使用历史。

教学掩码分数设正确记忆为0.8，参考为0.4，则$B=0.4$。同场景错误历史得0.7，恢复比例为$(0.7-0.4)/0.4=0.75$；其他对象历史得0.45，比例为0.125。前者说明大量收益相对这个对照保留下来，后者说明更换对象后多数收益消失。若某对照得0.35，恢复比例为$-0.125$，不应强制记为零；同时报告原分数才知道这种负比例发生在哪个质量水平。身份和事件输出需各自重做审计，不能用掩码分数的变化替它们判断依赖。

原论文在冻结视频模型的指定设置中发现，某些无实例内容对照仍保留较多收益；WorldMem对场景上下文和具体位置显示不同程度依赖，SAM 2的空间记忆替换则明显破坏目标分割，尤其在重新出现时。这里的“表示修复”解释还由离训练记忆的距离与固定幅度不同方向的对照支撑，不能从一次高恢复比例直接推出。模型、任务与消费接口各不相同，这组结果恰说明不能概括为所有记忆都忽略历史，也不能把SAM 2的内容依赖推广为所有视频系统的保证。

\section{本章小结}
\label{sec:vision-video-summary}

十帧中的同一辆车需要通过对应和状态管理归为一个身份，拿起与放下需要借助时间顺序和对象状态区分，找到某次动作则要保留查询所需的片段与边界。更多帧提供了补充证据，也增加了错配、压缩和等待的可能。

对象识别、轨迹维护、动态点查询、事件发现和检索各自保留输出职责，才能把错误定位到具体环节。局部表示要接入任务监督，三维运动要固定坐标与尺度，未完成事件还要按当时的可见帧判断响应。对每个结果说明实际看过哪些帧、哪些历史仍可访问、哪些位置由预测补出，是把局部模型组合成可核验视频应用的基础。回查的稳定答案仍需覆盖核验，记忆收益也需内容替换对照，才可判断正确性背后用了什么依据。

长期身份、细粒度活动、低预算观察和未完成事件响应仍是需要研究的问题。增加记忆容量只有在保留了任务可用信息时才有用，其中是否必须保留实例内容还需对照检验；缩短等待则可能减少可见的运动阶段。应把采样频率、有效历史长度、显存、计算时间和等待时间与任务分项结果一起报告，并检查被淘汰的观测会影响哪些未来查询。更值得继续检验的是怎样把局部状态、原始证据覆盖和及时输出放在同一预算协议下比较，而非只增加一个总分。\ref{chap:vision-video-restoration}章沿着可靠对应关系，讨论如何恢复和传输连续画面本身。
